\documentclass[11pt,a4paper]{article}

\usepackage[utf8]{inputenc}
\usepackage[spanish]{babel}
\usepackage{amsmath}
\usepackage{graphicx}
\usepackage{hyperref}
\usepackage{booktabs}
\usepackage{geometry}
\usepackage{caption}
\usepackage{float}
\usepackage{titlesec}
\usepackage{xcolor}

\geometry{margin=2.5cm}
\titleformat{\section}{\large\bfseries\color{darkblue}}{}{0em}{}
\definecolor{darkblue}{rgb}{0.1, 0.2, 0.5}

\title{
    \textbf{Sistema de Predicción de Churn Global-E-Shop}\\
    \large Informe Técnico de Arquitectura y Modelado E2E
}
\author{
    \textbf{Staff Data Engineer \& Lead Data Scientist}\\
    Departamento de Advanced Analytics
}
\date{\today}

\begin{document}

\maketitle

\begin{abstract}
Este documento detalla el diseño, implementación y validación del sistema de predicción de abandono de clientes (Churn) para Global-E-Shop. Se describe un pipeline end-to-end que integra una base de datos relacional, un pipeline de ingeniería de variables con rigor temporal, un modelo basado en XGBoost optimizado para Recall y una capa de serving mediante FastAPI. Los resultados demuestran un Lift de 1.87x en el segmento de mayor riesgo, permitiendo la automatización de campañas de retención con alta eficiencia operacional.
\end{abstract}

\section{Introducción}
La retención de clientes es el pilar de rentabilidad en el e-commerce. Identificar a los clientes en riesgo de abandono antes de que ocurra permite activar palancas de marketing personalizadas. El objetivo de este proyecto es construir un sistema robuso, reproducible y auditable que transforme datos transaccionales en decisiones de negocio accionables.

\section{Arquitectura del Sistema}

\subsection{Data Foundation (SQL)}
Se diseñó un esquema relacional normalizado que separa dimensiones (\texttt{dim\_customers}, \texttt{dim\_products}) de hechos (\texttt{fact\_orders}, \texttt{fact\_support\_tickets}). Esta estructura permite una integridad referencial estricta y facilita agregaciones temporales. Se incluyó un \texttt{ml\_feature\_store} para registrar las variables utilizadas en cada ciclo de entrenamiento, garantizando la auditabilidad del historial del modelo.

\subsection{Pipeline ETL y No-Leakage Temporal}
La ingeniería de variables se rige por la regla de la fecha de corte ($t_0$). Solo se utiliza información transaccional donde $t \le t_0$. El pipeline implementado realiza una limpieza de datos, manejo de nulos por imputación estadística y validación de invariantes de negocio (e.g., fechas de entrega coherentes).

\section{Metodología de Modelado}

\subsection{Ingeniería de Variables}
Se calcularon features de tres categorías principales:
\begin{itemize}
    \item \textbf{RFM}: Recency (días sin compra), Frequency (nro. pedidos) y Monetary (gasto total).
    \item \textbf{Intensity}: Gasto diario promedio en la ventana de observación.
    \item \textbf{Trend \& Diversity}: Pendiente de la tendencia de gasto y diversidad de categorías compradas.
\end{itemize}

\subsection{Algoritmo y Entrenamiento}
Se utilizó \textbf{XGBoost} debido a su superioridad en datos tabulares y manejo nativo de nulos. Para mitigar el desbalance de clases (~10\% de churn), se aplicó una estrategia combinada de:
\begin{enumerate}
    \item \textbf{SMOTE}: Para suavizar la frontera de decisión en el conjunto de entrenamiento.
    \item \textbf{Classification Weights}: Mediante el parámetro \texttt{scale\_pos\_weight} calculado dinámicamente.
\end{enumerate}

\section{Resultados y Evaluación}

\subsection{Métricas de Negocio}
Se priorizó el \textbf{Recall} (Sensibilidad) sobre el Accuracy, dado que el costo de un falso negativo (no detectar a un cliente que se va) es significativamente mayor al de un falso positivo. Utilizando un umbral optimizado de 0.2, se obtuvo un Recall en test de \textbf{0.578 [0.471, 0.685] (95\% CI)}.

\begin{table}[H]
\centering
\caption{Comparativa de Resultados}
\begin{tabular}{lcc}
\toprule
Métrica & Baseline (RFM) & Modelo XGBoost \\
\midrule
Recall (Churn) & 0.120 & \textbf{0.578} \\
Precision & 0.132 & 0.129 \\
AUC-ROC & 0.500 & \textbf{0.662} \\
Lift @ 20\% & 1.00 & \textbf{1.87} \\
\bottomrule
\end{tabular}
\end{table}

\subsection{Interpretatibilidad}
El análisis de importancia de variables reveló que la \textbf{Recency} (días desde la última compra) y la \textbf{Trend Slope} (caída del gasto) son los predictores primarios de abandono.

\section{Conclusiones}
El sistema implementado proporciona una base sólida para la automatización de la retención. Al capturar al 58\% de los churners potenciales contactando solo al 20\% de la población, el negocio puede optimizar su presupuesto de descuentos y mejorar el Customer Lifetime Value (CLV).

\nocite{*}
\bibliographystyle{plain}
\bibliography{references}

\end{document}
